\DocumentMetadata{
  pdfversion=2.0,
  pdfstandard=ua-2,
  lang=en-US,
  tagging=on,
  tagging-setup={math/setup=mathml-SE,tabsorder=structure}
}
\documentclass[11pt,letterpaper]{article}
\usepackage[margin=0.68in,includefoot]{geometry}
\usepackage{newcomputermodern}
\usepackage{amsmath,amscd,mathtools}
\usepackage{tikz,adjustbox}
\usetikzlibrary{arrows.meta,decorations.markings}
\providecommand{\square}{\mdlgwhtsquare}
\usepackage[hidelinks]{hyperref}
\hypersetup{
  pdftitle={Topology First-Year Exam, May 2022, Part 2},
  pdfauthor={University of Florida Department of Mathematics},
  pdfsubject={Graduate mathematics examination},
  pdfdisplaydoctitle=true,
  bookmarksopen=true
}
\setcounter{secnumdepth}{0}
\setlength{\parindent}{0pt}
\setlength{\parskip}{0.28em}
\setlength{\emergencystretch}{3em}
\setlength{\leftmargini}{2.15em}
\setlength{\leftmarginii}{2.65em}
\setlength{\leftmarginiii}{3em}
\renewcommand{\labelenumi}{\textbf{\theenumi.}}
\pagestyle{empty}
\raggedbottom
\allowdisplaybreaks
\newenvironment{examproblems}[1]{%
  \begin{enumerate}%
  \setcounter{enumi}{#1}%
  \setlength{\topsep}{0.35em}%
  \setlength{\partopsep}{0pt}%
  \setlength{\itemsep}{0.48em}%
  \setlength{\parsep}{0.18em}%
}{\end{enumerate}}
\newenvironment{examsubparts}{%
  \begin{enumerate}%
  \setlength{\topsep}{0.18em}%
  \setlength{\partopsep}{0pt}%
  \setlength{\itemsep}{0.16em}%
  \setlength{\parsep}{0.1em}%
}{\end{enumerate}}
\newenvironment{examinstructions}{%
  \par\begingroup\itshape
}{%
  \par\endgroup\vspace{0.18em}
}
\makeatletter
\renewcommand\section{\@startsection{section}{1}{0pt}%
  {-1.25ex plus -.25ex minus -.1ex}%
  {0.55ex plus .1ex}%
  {\normalfont\large\bfseries}}
\renewcommand{\@maketitle}{%
  \newpage\null\vskip -1em%
  {\footnotesize\bfseries
    \href{https://www.ufl.edu/}{University of Florida}, \href{https://math.ufl.edu/}{Department of Mathematics}\par}%
  \vskip -0.5em%
  \tagstructbegin{tag=Title}%
    {\LARGE\bfseries\href{https://gma.math.ufl.edu/exams/topology-first-year/}{Topology First-Year Exam}, May 2022, Part 2\par}%
  \tagstructend%
  \vskip 0.25em%
  \tagmcbegin{artifact}%
    \hrule height 1pt%
  \tagmcend%
  \vskip 1em%
}
\makeatother
\title{Topology First-Year Exam, May 2022, Part 2}
\author{}
\date{}
\begin{document}
\maketitle
\thispagestyle{empty}
\begin{examinstructions}
Do seven 7 and only seven of the following. In the following \(D^2\) is the closed, two-dimensional disk, and \(S^1\) is the unit circle, its boundary. The cyclic group of order~\(n\), is \(\mathbb{Z}_n=\mathbb{Z}/n\mathbb{Z}\).
\end{examinstructions}
\begin{examproblems}{0}
\item
Assume that~\(X\) is completely regular. Show that~\(X\) is connected if and only if \(\beta X\) is connected.
\item
This problem has three parts.
\begin{examsubparts}
\item[\textnormal{a.}]
Define a retract and show there is no retract \(D^2\to S^1\).
\item[\textnormal{b.}]
Show that any continuous map \(f:D^2\to D^2\) has a fixed point.
\item[\textnormal{c.}]
If~\(A\) is retract of \(D^2\) prove that any continuous map \(f:A\to A\) has a fixed point.
\end{examsubparts}
\item
Let~\(M\) and~\(N\) be two connected, compact manifolds with points \(x_0\in M\) and \(y_0\in N\). Define an equivalence relation on their disjoint union, \(M\sqcup N\), as \(x_0\sim y_0\) and all other points are just equivalent to themselves. Let \(M\vee N=(M\sqcup N)/{\sim}\). Show that 
\[
\pi_1(M\vee N,[x_0])=\pi_1(M,x_0)*\pi_1(N,y_0)
\]
 where \([x_0]\) is the equivalence class of \(x_0\).
\item
Construct a space whose fundamental group is isomorphic to \((\mathbb{Z}\oplus\mathbb{Z}_3)*\mathbb{Z}_5\).
\item
Show that \(S^2\) is not homeomorphic to \(S^n\) for any \(n>2\).
\item
State the Seifert–van Kampen Theorem and use it to compute the fundamental group of the space pictured below.
\par\smallskip
\begin{center}
\begin{adjustbox}{max width=.8\linewidth,max totalheight=6\baselineskip}
\begin{tikzpicture}[
  alt={A two-lobed surface crossing at its center. The left lobe contains a small oval hole, and the right lobe has a transverse circular cross-section whose hidden portion is dashed.},
  x=0.85cm,
  y=0.85cm,
  line width=1.15pt,
  line cap=round,
  line join=round
]
\coordinate (C) at (0,0);
\coordinate (P) at (2.90,2.10);
\coordinate (Q) at (3.45,-0.65);

% Left lobe
\draw
    (C)
    .. controls (-0.60,0.14) and (-1.75,1.12) .. (-3.35,1.45)
    .. controls (-4.35,1.66) and (-5.08,1.04) .. (-5.25,0.18)
    .. controls (-5.43,-0.72) and (-4.68,-1.66) .. (-3.55,-1.94)
    .. controls (-2.38,-2.22) and (-1.08,-0.88) .. (C);

% Right lobe
\draw
    (C)
    .. controls (0.66,0.30) and (1.84,1.55) .. (P)
    .. controls (3.80,2.57) and (4.57,2.43) .. (4.78,1.68)
    .. controls (5.05,0.70) and (4.39,-0.22) .. (Q)
    .. controls (2.42,-1.03) and (0.78,-0.40) .. (C);

% Small opening or handle in the left lobe
\draw
    (-4.12,-0.08)
    .. controls (-3.93,-0.36) and (-3.75,-0.57) .. (-3.52,-0.58)
    -- (-2.74,-0.58)
    .. controls (-2.52,-0.45) and (-2.34,-0.19) .. (-2.22,-0.05);

\draw
    (-3.52,-0.58)
    .. controls (-3.23,-0.20) and (-2.98,-0.19) .. (-2.74,-0.58);

% Visible half of the cross-section in the right lobe
\draw
    (P)
    .. controls (3.82,1.87) and (4.12,0.47) .. (Q);

% Hidden half of the cross-section
\draw[dash pattern=on 4.5pt off 4pt]
    (P)
    .. controls (2.66,1.03) and (2.83,0.05) .. (Q);
\end{tikzpicture}
\end{adjustbox}
\end{center}
\smallskip
\item
The projective plane \(PP\) is constructed by identifying the edges of a disk as pictured below. Compute the fundamental group of \(PP\) as well as that of \(PP\) minus~\(n\) points, \(n>0\).
\par\smallskip
\begin{center}
\begin{adjustbox}{max width=.8\linewidth,max totalheight=6\baselineskip}
\begin{tikzpicture}[
  alt={A lens-shaped disk with its upper and lower boundary arcs both labeled a. An arrow on the upper arc points right, and an arrow on the lower arc points left, indicating the edge identification.},
  line width=1.15pt,
  line cap=round,
  line join=round,
  boundary arrow/.style={
        postaction={
            decorate,
            decoration={
                markings,
                mark=at position 0.52 with
                    {\arrow{Latex[length=3.5mm,width=2.6mm]}}
            }
        }
    }
]
\coordinate (L) at (-4,0);
\coordinate (R) at ( 4,0);

% Upper edge, oriented from left to right
\draw[boundary arrow]
    (L)
    .. controls (-2.25,1.55) and (1.85,2.05) ..
    (R);

% Lower edge, oriented from right to left
\draw[boundary arrow]
    (R)
    .. controls (2.20,-1.55) and (-1.95,-1.85) ..
    (L);

% Edge labels
\node at (0,2.05)  {$a$};
\node at (0,-1.95) {$a$};

% Light interior hatching
\draw (-1.85,0.48) -- (-2.15,-0.20);
\draw (-0.95,0.55) .. controls (-1.18,0.15) .. (-1.20,-0.42);
\draw ( 0.00,0.45) -- (-0.25,-0.28);
\draw ( 0.90,0.36) .. controls (0.65,0.05) .. (0.70,-0.25);
\draw ( 1.75,0.30) -- (1.55,-0.42);
\end{tikzpicture}
\end{adjustbox}
\end{center}
\smallskip
\item
This problem has two parts.
\begin{examsubparts}
\item[\textnormal{(a)}]
Define a deformation retract and show it induces an isomorphism on fundamental groups.
\item[\textnormal{(b)}]
Prove or disprove: a retract induces an isomorphism on fundamental groups.
\end{examsubparts}
\end{examproblems}
\end{document}
